\section*{\texorpdfstring{\S\CurrentExerciseGroup}{第{\CurrentExerciseGroup}节}}
\phantomsection
\addcontentsline{toc}{section}{第{\CurrentExerciseGroup}节习题}

\begin{enumerate}
\def\labelenumi{\arabic{enumi})}
\tightlist
\item
  设 \(X, Y\) 是两个局部紧空间。
\end{enumerate}

\begin{enumerate}
\def\labelenumi{\alph{enumi})}
\item
  试证：映射 \((\mu, \nu) \mapsto \mu \otimes \nu\) （从 \(\mathcal{M}(\mathrm{X};\mathbf{C}) \times \mathcal{M}(\mathrm{Y};\mathbf{C})\) 到 \(\mathcal{M}(\mathrm{X} \times \mathrm{Y};\mathbf{C})\) 中）当 \(\mathcal{M}(\mathrm{X};\mathbf{C})\) 、\(\mathcal{M}(\mathrm{Y};\mathbf{C})\) 及 \(\mathcal{M}(\mathrm{X} \times \mathrm{Y};\mathbf{C})\) 均赋以拟强拓扑（§1，习题 8）时是连续的。
\item
  试证：映射 \((\mu, \nu) \mapsto \mu \otimes \nu\) （从 \(\mathcal{M}_{+}(X) \times \mathcal{M}_{+}(Y)\) 到 \(\mathcal{M}_{+}(X \times Y)\) 中）当 \(\mathcal{M}_{+}(X)\) 、\(\mathcal{M}_{+}(Y)\) 及 \(\mathcal{M}_{+}(X \times Y)\) 均赋以淡拓扑时是连续的。
\end{enumerate}

¶ 2) a) 设 X 是 R 的区间 {[}0,1{]} ；试证：若 \((a_k)_{1 \leqslant k \leqslant n}\) 是 X 中两两不同元素的任意有限序列，则 \(n\) 个函数 \(|x - a_k|\) （ \(1 \leqslant k \leqslant n\) ）线性无关。由此推出：连续函数 \(|x - y|\) （定义在 \(X \times X\) 上）不是 \(\sum_{i=1}^{n} u_i(x)v_i(y)\) 的形式，其中诸 \(u_i\) 与 \(v_i\) 连续。

\begin{enumerate}
\def\labelenumi{\alph{enumi})}
\setcounter{enumi}{1}
\item
  对每个测度 \(\mu \in \mathcal{M}(\mathrm{X};\mathbf{C})\) ，令 \(\mathbf{f}_{\mu}(y) = \int |x - y| d\mu(x)\) 。试证：线性映射 \(\mu \mapsto \mathbf{f}_{\mu}\) （从 \(\mathcal{M}(\mathrm{X};\mathbf{C})\) 到 \(\mathcal{C}(\mathrm{X};\mathbf{C})\) 中）是单射，其像是 \(\mathcal{C}(\mathrm{X};\mathbf{C})\) 的一个稠密线性子空间 L ，且对于 \(\mathcal{M}(\mathrm{X};\mathbf{C})\) 与 L 上的赋范空间拓扑连续但非双连续。（注意 L 包含常函数与分段线性函数。）
\item
  设 \(u_{i}\) ( \(1 \leqslant i \leqslant m\) ) 是任意的 \(m\) 个函数（属于 \(\mathcal{C}(\mathrm{X};\mathbf{C})\) ），设 \(V\) 是 \(\mathcal{M}(\mathrm{X};\mathbf{C})\) 中与诸 \(u_{i}\) 正交的线性子空间，因而其余维数 \(\leqslant m\) （在 \(\mathcal{M}(\mathrm{X};\mathbf{C})\) 中）。试证：若 \(B\) 是测度 \(\mu\) （在 \(X\) 上）所成的满足 \(\| \mu \| \leqslant 1\) 的集，则存在测度 \(\mu \in B\) 与 \(\nu \in V\) 使 \(\int \mathfrak{f}_{\mu}(y)d\nu(y)\) 可任意大（利用 \(b\) ）。由此推出：若 \(B\) 、\(\mathcal{M}(X; \mathbf{C})\) 及 \(\mathcal{M}(X \times X; \mathbf{C})\) 均赋以淡拓扑，则映射 \((\mu, \nu) \mapsto \mu \otimes \nu\) （从 \(B \times \mathcal{M}(X; \mathbf{C})\) 到 \(\mathcal{M}(X \times X; \mathbf{C})\) 中）不连续。
\end{enumerate}

¶ 3) a) 设 X 、Y 是两个局部紧空间，K ⊂ X 、L ⊂ Y 是两个紧集，A 是空间 \(\mathcal{K}(X \times Y, K \times L; \mathbf{C})\) 的一个紧子集。对每个整数 \(n\) ，设 \(V_n\) 是 K 的一致结构的一个近域，使得对每对 \((x', x'') \in V_n\) 、每个 \(y \in L\) 及每个函数 \(f \in A\) ，

\[
\left| f (x ^ {\prime}, y) - f (x ^ {\prime \prime}, y) \right| \leqslant 1 / n ^ {2}.
\]

设 \(C_n\) 是函数 \(y \mapsto n(f(x', y) - f(x'', y))\) （对所有对 \((x', x'') \in V_n\) 及每个 \(f \in A\) ）所成的集；试证：并 \(C\) ，即诸 \(C_n\) （对所有 \(n \geqslant 1\) ）的并，是 \(\mathcal{K}(Y, L; \mathbf{C})\) 的相对紧子集。

\begin{enumerate}
\def\labelenumi{\alph{enumi})}
\setcounter{enumi}{1}
\tightlist
\item
  设 B 是 C 与映射 \(y \mapsto f(x, y)\) （对 \(x \in K\) 及 \(f \in A\) ）所成之集的并，它是 \(\mathcal{K}(\mathrm{Y}, \mathrm{L}; \mathbf{C})\) 的一个相对紧子集。试证：当 \(\nu\) 遍历极 \(B^{\circ}\) （B 在 \(\mathcal{M}(\mathrm{Y}; \mathbf{C})\) 中的）且 f 遍历 A 时，映射所成之集
\end{enumerate}

\[
y \mapsto \int f (x, y) d \nu (y)
\]

在 \(\mathcal{H}(\mathrm{X},\mathrm{K};\mathbf{C})\) 中相对紧

\begin{enumerate}
\def\labelenumi{\alph{enumi})}
\setcounter{enumi}{2}
\tightlist
\item
  由 b) 推出：映射 \((\mu, \nu) \mapsto \mu \otimes \nu\) （从 \(\mathcal{M}(\mathrm{X}; \mathbf{C}) \times \mathcal{M}(\mathrm{Y}; \mathbf{C})\) 到 \(\mathcal{M}(\mathrm{X} \times \mathrm{Y}; \mathbf{C})\) 中）当 \(\mathcal{M}(\mathrm{X}; \mathbf{C})\) 、\(\mathcal{M}(\mathrm{Y}; \mathbf{C})\) 及 \(\mathcal{M}(\mathrm{X} \times \mathrm{Y}; \mathbf{C})\) 均赋以严格紧收敛拓扑时是连续的。
\end{enumerate}

\begin{enumerate}
\def\labelenumi{\arabic{enumi})}
\setcounter{enumi}{3}
\item
  设 X 是局部紧空间，Y 是仿紧的局部紧空间。试证：\(\mathcal{K}(X \times Y; \mathbf{C})\) 到 \(\mathcal{K}(X; \mathcal{K}(Y; \mathbf{C}))\) 中的典范映射是拓扑向量空间的一个同构。
\item
  设 \((\mathrm{X}_{\lambda})_{\lambda \in L}\) 是紧空间的一个无穷族；对每个 \(\lambda \in L\) ，设 \(a_{\lambda}\) 是 \(\mathrm{X}_{\lambda}\) 的一个点，\(\mu_{\lambda}\) 是 \(\mathrm{X}_{\lambda}\) 上总质量为 1 的正测度。对每个有限子集 \(J\) ⊂ \(L\) ，设 \(\nu(J)\) 是 \(X = \prod_{\lambda \in L} X_{\lambda}\) 上的这样的测度，它是诸
\end{enumerate}

测度 \(\mu_{\lambda}\) （对 \(\lambda \in J\) ）与诸测度 \(\varepsilon_{a_{\lambda}}\) （对 \(\lambda \in L - J\) ）的乘积。试证：关于 \(L\) 的有限子集组成的有向集，测度 \(\nu(J)\) 淡收敛于 \(\mu = \bigotimes_{\lambda \in L} \mu_{\lambda}\)

但不强收敛于 \(\mu\) 。

¶ 6) 采用 No.~6 的记号，试证：为使 X 上存在测度 \(\mu \neq 0\) 满足

\[
\mu \left(f _ {J} \circ \operatorname{pr} _ {J}\right) = \lim _ {K \supset J} \mu_ {K} \left(f _ {J} \circ \operatorname{pr} _ {J, K}\right) = \mu_ {J} \left(f _ {J}\right) \cdot \prod_ {\lambda \in L - J} \mu_ {\lambda} (1)\tag{1}
\]

对每个有限子集 \(\mathbf{J}\) ⊂ \(\mathbf{L}\) 及每个函数 \(f_{\mathbf{J}} \in \mathcal{C}(\mathrm{X}_{\mathbf{J}}; \mathbf{C})\) ，必须且只须下述三个条件成立：

\(1^{\circ} \mu_{\lambda} \neq 0\) 对所有 \(\lambda \in L\) 成立。

\(2^{\circ}\) 存在有限子集 \(\mathbf{J}_0\) ⊂ \(\mathbf{L}\) ，使族 \((\mu_{\lambda}(1))_{\lambda \in \mathbf{L} - \mathbf{J}_0}\) 在 \(\mathbf{C}^*\) 中可乘（从而有乘积 \(\neq 0\) ）。

\(3^{\circ}\) 族 \((\| \mu_{\lambda}\|)_{\lambda \in L}\) 在 \(\mathbf{R}_{+}^{*}\) 中可乘（从而有乘积 \(\neq 0\) ）。

在什么条件下，(1) 的右端对 L 的每个有限子集 J 及每个函数 \(f_{\mathrm{J}} \in \mathcal{C}(\mathrm{X}_{\mathrm{J}}; \mathbf{C})\) 都等于 0 ？

\begin{enumerate}
\def\labelenumi{\arabic{enumi})}
\setcounter{enumi}{6}
\tightlist
\item
  采用习题 6 的记号，并假设该习题的条件成立。试证：对每个函数 \(f \in \mathcal{C}(\mathrm{X};\mathbf{C})\) ，
\end{enumerate}

\[
\int f d \mu = \lim _ {\mathrm{J}} \int f (x _ {\mathrm{J}}, x _ {\mathrm{L} - \mathrm{J}}) d \mu_ {\mathrm{J}} (x _ {\mathrm{J}}),
\]

其中极限是关于 L 的有限子集组成的有向集取的，且 \(x_{J}\) 表示 \(\Pr_{J}(x)\) （对每个 \(x \in X\) 及每个 \(J \subset L\) ）；类似地，

\[
f (x) = \lim _ {\mathrm{J}} \int f (x _ {\mathrm{J}}, x _ {\mathrm{L} - \mathrm{J}}) d \mu_ {\mathrm{L} - \mathrm{J}} (x _ {\mathrm{L} - \mathrm{J}})
\]

对每个 \(x \in X\) 成立，只要诸测度 \(\mu_{\lambda} (\lambda \in L)\) 都具有总质量 1 （这里，对每个子集 \(K\) ⊂ \(L\) ，\(\mu_{K}\) 表示测度子族 \((\mu_{\lambda})_{\lambda \in K}\) 的乘积）。

\begin{enumerate}
\def\labelenumi{\arabic{enumi})}
\setcounter{enumi}{7}
\tightlist
\item
  采用习题 6 的记号，试证：若族 \((\mu_{\lambda})\) 的乘积有定义且 \(\neq 0\) ，则存在可数子集 \(K\) ⊂ \(L\) ，使对每个 \(\lambda \in L - K\) ，测度 \(\mu_{\lambda}\) 是正的且总质量为 1 （利用 §1 的习题 18）。
\end{enumerate}

